% GENERATED FILE. Edit canonical lessons, then run npm run generate:books.
% canonical-source-hash: fnv1a64:a809ff206b653aed
% canonical-lessons: KA-C422-paramarsisu, KA-C422-manadanda, KA-C422-krama, KA-C422-vidhana, KA-C422-prakriye

\chapter{Rules and Methods}
\label{ch:ka-rules-and-methods}
\hlchaptermodality{422}

\begin{chapteropening}
\textbf{By the end of this chapter:} \emph{I can say to refer to, a criterion, a course of action, a method and a process.}

Complete the last lesson of chapter 422: I can say to refer to, a criterion, a course of action, a method and a process.
\end{chapteropening}

\section[\texorpdfstring{\kn{ಪರಾಮರ್ಶಿಸು} (parāmarśisu)}{ಪರಾಮರ್ಶಿಸು (parāmarśisu)}]{\kn{ಪರಾಮರ್ಶಿಸು} (parāmarśisu) — to refer to; to consult}

\label{lesson:KA-C422-paramarsisu}



\begin{quote}
Before the new one: say the Kannada for export, then the Kannada for inflation.
\end{quote}

\subsection*{You'll want to know: \kn{ಪರಾಮರ್ಶಿಸು}}
\textbf{\kn{ಪರಾಮರ್ಶಿಸು}} — \emph{parāmarśisu} — \textquotedblleft{}to refer to; to consult\textquotedblright{}.

\begin{quote}
\kn{ಉತ್ತರ} \kn{ಬರೆಯುವ} \kn{ಮೊದಲು} \kn{ನಾನು} \kn{ನಿಘಂಟನ್ನು} \kn{ಪರಾಮರ್ಶಿಸಿದೆ}.
\end{quote}

\begin{quote}
\emph{uttara bareyuva modalu nānu nighaṇṭannu parāmarśiside.}
\end{quote}

\begin{quote}
Before writing the answer I consulted the dictionary.
\end{quote}

Past (I …): \kn{ಪರಾಮರ್ಶಿಸಿದೆ}, \emph{parāmarśiside}.

Of books and records. \kn{ಪರಾಮರ್ಶನ} \kn{ಗ್ರಂಥ}, a reference book.

\begin{cousinweb}
\textbf{\kn{ಪರಾಮರ್ಶಿಸು}} (\emph{parāmarśisu}) is built on Sanskrit \emph{parāmarśa}, \textquotedblleft{}seizing; consideration\textquotedblright{} (Monier-Williams), with the Kannada verb ending -\emph{isu}.
\end{cousinweb}

\begin{grammarlens}[title={What you've built}]
One more word for saying what happens if.
\end{grammarlens}

\subsection*{Your turn}
\begin{itemize}
\raggedright
  \item \emph{Say it:} \emph{parāmarśisu}
  \item \emph{Say it:} \emph{parāmarśisu}, once more
  \item \emph{Say it:} \emph{uttara bareyuva modalu nānu nighaṇṭannu parāmarśiside.} — Before writing the answer I consulted the dictionary.
  \item \emph{Recall:} say the Kannada for export, then the Kannada for inflation, then say \emph{parāmarśisu} again
\end{itemize}

\begin{tcolorbox}[breakable,colback=teal!4,colframe=teal!35!black,title={Before you move on}]
What does \kn{ಪರಾಮರ್ಶಿಸು} mean? (\textquotedblleft{}To refer to; to consult\textquotedblright{}.) Say it once more.
\end{tcolorbox}
\section[\texorpdfstring{\kn{ಮಾನದಂಡ} (mānadaṇḍa)}{ಮಾನದಂಡ (mānadaṇḍa)}]{\kn{ಮಾನದಂಡ} (mānadaṇḍa) — a criterion; a standard}

\label{lesson:KA-C422-manadanda}



\begin{quote}
Before the new one: say the Kannada for inflation, then the Kannada for to refer to.
\end{quote}

\subsection*{You'll want to know: \kn{ಮಾನದಂಡ}}
\textbf{\kn{ಮಾನದಂಡ}} — \emph{mānadaṇḍa} — \textquotedblleft{}a criterion; a standard\textquotedblright{}.

\begin{quote}
\kn{ಆಯ್ಕೆಗೆ} \kn{ಅಂಕಗಳೇ} \kn{ಮುಖ್ಯ} \kn{ಮಾನದಂಡ}.
\end{quote}

\begin{quote}
\emph{āykege aṅkagaḷē mukhya mānadaṇḍa.}
\end{quote}

\begin{quote}
Marks are the main criterion for selection.
\end{quote}

\kn{ಮಾನ} (measure) + \kn{ದಂಡ} (rod): a measuring rod.

\begin{cousinweb}
\textbf{\kn{ಮಾನದಂಡ}} (\emph{mānadaṇḍa}) is built from Sanskrit \emph{māna}, \textquotedblleft{}a measure\textquotedblright{} (Monier-Williams), and \emph{daṇḍa}, \textquotedblleft{}a stick, a rod\textquotedblright{} (CDIAL 6128): a measuring rod.
\end{cousinweb}

\begin{grammarlens}[title={What you've built}]
One more word for saying what happens if.
\end{grammarlens}

\subsection*{Your turn}
\begin{itemize}
\raggedright
  \item \emph{Say it:} \emph{mānadaṇḍa}
  \item \emph{Say it:} \emph{mānadaṇḍa}, once more
  \item \emph{Say it:} \emph{āykege aṅkagaḷē mukhya mānadaṇḍa.} — Marks are the main criterion for selection.
  \item \emph{Recall:} say the Kannada for inflation, then the Kannada for to refer to, then say \emph{mānadaṇḍa} again
\end{itemize}

\begin{tcolorbox}[breakable,colback=teal!4,colframe=teal!35!black,title={Before you move on}]
What does \kn{ಮಾನದಂಡ} mean? (\textquotedblleft{}A criterion; a standard\textquotedblright{}.) Say it once more.
\end{tcolorbox}
\section[\texorpdfstring{\kn{ಕ್ರಮ} (krama)}{ಕ್ರಮ (krama)}]{\kn{ಕ್ರಮ} (krama) — a course of action; a measure}

\label{lesson:KA-C422-krama}



\begin{quote}
Before the new one: say the Kannada for to refer to, then the Kannada for a criterion.
\end{quote}

\subsection*{You'll want to know: \kn{ಕ್ರಮ}}
\textbf{\kn{ಕ್ರಮ}} — \emph{krama} — \textquotedblleft{}a course of action; a measure\textquotedblright{}.

\begin{quote}
\kn{ಕಳ್ಳತನ} \kn{ತಡೆಯಲು} \kn{ಪೊಲೀಸರು} \kn{ಹೊಸ} \kn{ಕ್ರಮ} \kn{ತೆಗೆದುಕೊಂಡರು}.
\end{quote}

\begin{quote}
\emph{kaḷḷatana taḍeyalu polīsaru hosa krama tegedukoṇḍaru.}
\end{quote}

\begin{quote}
The police took new measures to stop theft.
\end{quote}

\kn{ಕ್ರಮವಾಗಿ}, in order; \kn{ಕ್ರಮ} \kn{ಕೈಗೊಳ್ಳು}, to take action.

\begin{cousinweb}
\textbf{\kn{ಕ್ರಮ}} (\emph{krama}) is borrowed from Sanskrit \emph{krama}, \textquotedblleft{}a step, going, order\textquotedblright{} (CDIAL 3577).
\end{cousinweb}

\begin{grammarlens}[title={What you've built}]
One more word for saying what happens if.
\end{grammarlens}

\subsection*{Your turn}
\begin{itemize}
\raggedright
  \item \emph{Say it:} \emph{krama}
  \item \emph{Say it:} \emph{krama}, once more
  \item \emph{Say it:} \emph{kaḷḷatana taḍeyalu polīsaru hosa krama tegedukoṇḍaru.} — The police took new measures to stop theft.
  \item \emph{Recall:} say the Kannada for to refer to, then the Kannada for a criterion, then say \emph{krama} again
\end{itemize}

\begin{tcolorbox}[breakable,colback=teal!4,colframe=teal!35!black,title={Before you move on}]
What does \kn{ಕ್ರಮ} mean? (\textquotedblleft{}A course of action; a measure\textquotedblright{}.) Say it once more.
\end{tcolorbox}
\section[\texorpdfstring{\kn{ವಿಧಾನ} (vidhāna)}{ವಿಧಾನ (vidhāna)}]{\kn{ವಿಧಾನ} (vidhāna) — a method; a procedure}

\label{lesson:KA-C422-vidhana}



\begin{quote}
Before the new one: say the Kannada for a criterion, then the Kannada for a course of action.
\end{quote}

\subsection*{You'll want to know: \kn{ವಿಧಾನ}}
\textbf{\kn{ವಿಧಾನ}} — \emph{vidhāna} — \textquotedblleft{}a method; a procedure\textquotedblright{}.

\begin{quote}
\kn{ಅಮ್ಮನ} \kn{ಅಡುಗೆಯ} \kn{ವಿಧಾನ} \kn{ಬೇರೆ}.
\end{quote}

\begin{quote}
\emph{ammana aḍugeya vidhāna bēre.}
\end{quote}

\begin{quote}
Mother's way of cooking is different.
\end{quote}

\kn{ವಿಧಾನಸಭೆ}, the legislative assembly.

\begin{cousinweb}
\textbf{\kn{ವಿಧಾನ}} (\emph{vidhāna}) is borrowed from Sanskrit \emph{vidhāna}, \textquotedblleft{}arranging, regulating, a rule\textquotedblright{} (CDIAL 11754).
\end{cousinweb}

\begin{grammarlens}[title={What you've built}]
One more word for saying what happens if.
\end{grammarlens}

\subsection*{Your turn}
\begin{itemize}
\raggedright
  \item \emph{Say it:} \emph{vidhāna}
  \item \emph{Say it:} \emph{vidhāna}, once more
  \item \emph{Say it:} \emph{ammana aḍugeya vidhāna bēre.} — Mother's way of cooking is different.
  \item \emph{Recall:} say the Kannada for a criterion, then the Kannada for a course of action, then say \emph{vidhāna} again
\end{itemize}

\begin{tcolorbox}[breakable,colback=teal!4,colframe=teal!35!black,title={Before you move on}]
What does \kn{ವಿಧಾನ} mean? (\textquotedblleft{}A method; a procedure\textquotedblright{}.) Say it once more.
\end{tcolorbox}
\section[\texorpdfstring{\kn{ಪ್ರಕ್ರಿಯೆ} (prakriye)}{ಪ್ರಕ್ರಿಯೆ (prakriye)}]{\kn{ಪ್ರಕ್ರಿಯೆ} (prakriye) — a process}

\label{lesson:KA-C422-prakriye}



\begin{quote}
Before the new one: say the Kannada for a course of action, then the Kannada for a method.
\end{quote}

\subsection*{You'll want to know: \kn{ಪ್ರಕ್ರಿಯೆ}}
\textbf{\kn{ಪ್ರಕ್ರಿಯೆ}} — \emph{prakriye} — \textquotedblleft{}a process\textquotedblright{}.

\begin{quote}
\kn{ಮನೆ} \kn{ಕೊಳ್ಳುವ} \kn{ಪ್ರಕ್ರಿಯೆ} \kn{ತುಂಬಾ} \kn{ಉದ್ದ}.
\end{quote}

\begin{quote}
\emph{mane koḷḷuva prakriye tumbā udda.}
\end{quote}

\begin{quote}
The process of buying a house is very long.
\end{quote}

\kn{ಪ್ರಕ್ರಿಯೆ} \kn{ಮುಗಿಸು}, to complete a process.

\begin{cousinweb}
\textbf{\kn{ಪ್ರಕ್ರಿಯೆ}} (\emph{prakriye}) is borrowed from Sanskrit \emph{prakriyā}, \textquotedblleft{}production; procedure, way, manner\textquotedblright{} (Monier-Williams).
\end{cousinweb}

\begin{grammarlens}[title={What you've built}]
One more word for saying what happens if.
\end{grammarlens}

\subsection*{Your turn}
\begin{itemize}
\raggedright
  \item \emph{Say it:} \emph{prakriye}
  \item \emph{Say it:} \emph{prakriye}, once more
  \item \emph{Say it:} \emph{mane koḷḷuva prakriye tumbā udda.} — The process of buying a house is very long.
  \item \emph{Recall:} say the Kannada for a course of action, then the Kannada for a method, then say \emph{prakriye} again
\end{itemize}

\begin{tcolorbox}[breakable,colback=teal!4,colframe=teal!35!black,title={Before you move on}]
What does \kn{ಪ್ರಕ್ರಿಯೆ} mean? (\textquotedblleft{}A process\textquotedblright{}.) Say it once more.
\end{tcolorbox}
